% ─────────────────────────────────────────────────────────────
\phantomsection
\chapter{Introduction}
\label{sec:introduction}

% ─────────────────────────────────────────────────────────────
\section{Clinical and Societal Motivation}
\label{subsec:motivation}

Falls are the second leading cause of accidental injury death worldwide, after road
traffic injuries~\parencite{who2021falls}.  Globally, approximately 684,000 fatal falls
occur each year, with adults over the age of 60 accounting for the overwhelming
majority~\parencite{who2021falls}.  An estimated 37.3~million falls annually are severe
enough to require medical attention, representing a sustained and growing burden on
healthcare systems.  As the global population ages---the World Health Organization
projects that the number of people aged 60 and above will reach 2.1~billion by
2050~\parencite{who2022ageing}---the scale of this burden is set to increase substantially
without effective countermeasures.

The medical consequences of a fall extend well beyond the acute injury.  A fall that
goes undetected initiates a ``long lie'': a period during which the fallen person
remains on the ground, unable to summon help, and is exposed to secondary complications
including dehydration, hypothermia, pressure sores, and aspiration pneumonia.  Even
when the acute injury is minor, a long lie lasting more than an hour doubles the
probability of subsequent institutionalisation and is independently associated with
increased one-year mortality~\parencite{wild1981longlie}.  Beyond the physical consequences,
the \emph{fear of falling}---a well-documented post-fall psychological
syndrome~\parencite{tinetti1988fear}---leads to self-imposed activity restriction that
accelerates deconditioning and, paradoxically, further elevates fall risk.  Timely
detection is therefore not merely about comfort; it is a clinical intervention with
measurable outcomes.

Wearable inertial sensors offer a continuous and unobtrusive monitoring modality that
is uniquely positioned to close this detection gap.  A single accelerometer attached
to the waist or carried in a trouser pocket records the three-axis acceleration of the
body's centre of mass throughout the day.  Unlike camera-based or floor-sensor systems,
it imposes no infrastructure requirement and preserves the user's privacy and mobility.
Modern smartphones and smartwatches already contain the necessary hardware and are
carried by a large fraction of the target population, making passive fall detection a
feature that can be deployed at scale at marginal hardware cost.

The contrast between the two kinds of error is the premise that shapes the rest of
this thesis. A missed fall can trigger the long-lie complications described above,
whereas a false alarm costs at most a brief, unnecessary check; the two outcomes are
therefore far from equally serious. For this reason the present work pursues a
\emph{recall-oriented} objective and scores every model with the $F_2$ measure, which
weights recall four times as heavily as precision. In effect the detector is
deliberately tuned to give precedence to the fall event---tolerating the occasional
false alarm rather than risk overlooking a genuine fall---which is the design stance
reflected in the title of this thesis. The formal definition of $F_2$ and the reasoning
behind the four-to-one weighting are developed in the Methodology.

% ─────────────────────────────────────────────────────────────
\section{Limitations of Existing Approaches}
\label{subsec:limitations-related}

Despite the practical appeal of wearable fall detection, the path from sensor signal
to reliable deployment is obstructed by a set of well-known technical difficulties.
A brief taxonomy of existing approaches clarifies where the gaps lie.

\paragraph{Threshold-based detectors.}
The earliest and most widely deployed family of methods compares the signal magnitude
vector $\|\mathbf{a}(t)\|_2$ against a fixed threshold~$T$: a fall is declared
whenever $\|\mathbf{a}(t)\|_2 > T$.  These methods are computationally trivial and
require no training data.  Their fundamental limitation is the absence of a principled
calibration mechanism.  The threshold~$T$ is set empirically, the output is binary
with no associated uncertainty, and there is no statistical framework for adjusting
the operating point to the false-alarm tolerance of a particular deployment environment.
In practice, fixed-threshold detectors achieve high recall on controlled datasets but
at false-alarm rates that are clinically unacceptable for continuous wear: vigorous
activities of daily living (running, stair descent, jumping) frequently exceed the
same magnitude threshold as a fall impact.

\paragraph{Statistical and spectral feature methods.}
A substantial literature constructs handcrafted feature vectors from short
accelerometer windows---means, variances, spectral energies, signal-magnitude
area, zero-crossing rates---and trains standard classifiers on these
features~\parencite{igual2013challenges, delahoz2014survey}.  These methods improve over
threshold-based approaches by producing calibrated probability outputs and by
exploiting the multi-dimensional structure of the signal beyond a single scalar
statistic.  However, their features are \emph{ad hoc}: selected by domain intuition
rather than derived from a principled geometric theory of the signal.  There is
no formal stability guarantee---small perturbations of the input signal can cause
arbitrary changes in the feature vector if the perturbation happens to cross a
feature threshold.  As a result, the transferability of trained models across
sensors, sampling rates, and subject populations is empirically fragile.

\paragraph{Deep learning approaches.}
Convolutional and recurrent neural networks have achieved competitive performance on
fall detection benchmarks by learning end-to-end feature representations from raw
inertial data~\parencite{batchuluun2022deep, wang2020falldl}.  The gains over handcrafted
features are real, but the engineering trade-offs are severe.  Training requires large
labelled datasets, GPU compute, and hyperparameter searches whose cost scales with
model complexity.  The resulting models are opaque: the learned features are not
interpretable, and there is no geometric or mathematical explanation for why a
particular signal window is classified as a fall.  Deployment on edge hardware
(smartwatches, embedded sensors) is complicated by memory footprint and energy
consumption.  Most critically for the personalization question, deep models are
trained once on a fixed corpus and deployed without adaptation---the ``trains once,
deploys forever'' paradigm---which systematically ignores the substantial variation
in motion dynamics across individuals.

\paragraph{The personalization gap.}
Across all three families, a common structural weakness emerges: the decision rule is
calibrated at the population level and applied uniformly to every individual.  This
is statistically equivalent to assuming that the optimal decision boundary is the
same for all subjects.  In reality, individuals differ in gait dynamics, body geometry,
sensor placement geometry, and fall style, so the optimal threshold for a classifier
applied to subject~$s$ will in general differ from the population-level optimum.
Recognising this, some recent works report \emph{personalized} evaluation results, but
typically through isolated case studies rather than a systematic statistical
characterisation.  The question of whether subject-level adaptation is uniformly
beneficial, or only selectively so, remains open.

% ─────────────────────────────────────────────────────────────
\section{Topological Data Analysis for Wearable Signals}
\label{subsec:tda-motivation}

Persistent homology, the central tool of topological data analysis
(TDA)~\parencite{edelsbrunner2010computational, carlsson2009topology}, provides an
alternative geometric framework for characterising the shape of a signal window.
Rather than extracting fixed statistics from a time series, it constructs a sequence
of simplicial complexes---the \emph{filtration}---that tracks how the topology of the
point cloud derived from the signal evolves across spatial scales.  The result is a
\emph{persistence diagram}: a multiset of $(b, d)$ pairs recording the scales at
which topological features (connected components, loops, voids) are born and die.
This diagram summarises the intrinsic geometry of the signal window without
committing to a fixed set of scalar statistics.

Two properties of persistent homology are especially relevant to the fall-detection
problem.

\textbf{Provable stability.}  The Cohen--Steiner stability theorem~\parencite{cohen2007stability}
states that the bottleneck distance between two persistence diagrams is bounded above
by the supremum norm of the perturbation between the corresponding input functions:
\[
  d_B(D(f),\,D(g)) \;\leq\; \|f - g\|_\infty.
\]
Applied to inertial windows, this guarantees that small sensor noise perturbations
produce at most equally small perturbations of the topological descriptor.  This
stability is a \emph{mathematical theorem}, not an empirical observation, and it holds
independently of the training dataset and subject population.  No existing handcrafted
feature family possesses a comparable guarantee.

\textbf{Fixed-dimensional vectorization.}  The persistence diagram is converted into
a fixed 24-dimensional feature vector $\phi_\lambda(w)\in\R^{24}$ by computing
summary statistics of the birth-death pairs in homological degrees~$k=0$ and~$k=1$.
The dimension 24 is independent of the window length, sampling rate, or number of
simplices in the filtration.  This is in contrast to spectral or autoregressive
features whose dimension grows with the signal length, and to deep-network embeddings
whose dimension is a design parameter affecting model size.

The \emph{delay embedding}~\parencite{takens1981detecting} is the geometric pre-processing
step that connects the scalar signal to the point cloud: the orientation-invariant
signal magnitude vector $\mathrm{SMV}(t)=\|\mathbf{a}(t)\|_2$ is lifted to a
$m$-dimensional point cloud by taking overlapping $m$-tuples of lagged values.  A
fall window and an ADL window produce point clouds of qualitatively different
topology: the sharp impulsive excursion of a fall impact generates a spread-out
cluster with strong $H_0$ structure, whereas ADL windows exhibit richer $H_1$
structure (loops) from the periodic motion of walking or stair climbing.  Persistent
homology captures this topological distinction systematically.

TDA has been applied to time-series classification problems across several
domains~\parencite{gidea2018topological, pereira2015persistent, umeda2017time}, and
specifically to activity recognition from inertial
sensors~\parencite{xu2019topological, ott2017topological}.  The distinguishing features
of the present work are its recall-oriented objective, its systematic dataset-specific
hyperparameter optimisation, and its statistical treatment of subject-level
personalisation as a distributional question rather than a case study.

% ─────────────────────────────────────────────────────────────
\section{The Threshold Decision as a First-Class Component}
\label{subsec:threshold-intro}

Wearable fall detection is an instance of binary classification under class imbalance:
falls are rare relative to activities of daily living, so the positive class is
significantly underrepresented in any realistic deployment.  Under such imbalance,
the choice of decision threshold $\tau$ in the rule
$\widehat{y}(w)=\mathbf{1}\{f_\theta(\phi_\lambda(w))\geq\tau\}$
is not a cosmetic detail but a critical design parameter that directly controls the
precision-recall trade-off.

As established in Section~\ref{subsec:motivation}, the thesis prioritises \emph{recall}
through the recall-weighted $F_2$ score. Under class imbalance the decision threshold
$\tau$ is the concrete lever that realises this preference: lowering $\tau$ raises
recall at the expense of precision, so the operating point must be chosen deliberately
rather than left at the default $0.5$.

Beyond the population-level threshold, the thesis investigates \emph{subject-level
threshold personalisation}.  The key observation is that the posterior probability
function $f_\theta\circ\phi_{\lambda^*}$ may be systematically miscalibrated
differently for different subjects: a subject whose fall windows produce
lower-than-average posterior scores needs a lower threshold to achieve high recall,
while a subject whose ADL windows produce elevated scores needs a higher threshold to
control false alarms.  A sweep of $\tau$ over the grid
$\mathcal{T}=\{0.05,0.10,\ldots,0.95\}$ identifies the per-subject optimal threshold
$\tau^*(s)$ retrospectively.

The computational cost of this sweep is $O(|\mathcal{T}|\cdot|W_s|)$ comparisons,
where $|\mathcal{T}|=19$ and $|W_s|$ is the window count for subject~$s$.  This
amounts to less than one second of wall time for all subjects in any of the three
datasets, against a Bayesian representation search that consumes approximately
75~minutes per dataset.  The ratio exceeds $10^3$: threshold personalisation is a
computationally negligible add-on to an already-trained model.  It requires no model
retraining, no access to the original training data, and no transmission of data to
a server---it is executable on-device at inference time from a small per-subject
calibration set.  This combination of properties---near-zero cost, privacy
preservation, and measurable $F_2$ gain---makes the threshold decision a first-class
engineering contribution of the thesis, not merely an evaluation detail.

% ─────────────────────────────────────────────────────────────
\section{Contributions}
\label{subsec:contributions}

The thesis makes the following five contributions.

\textbf{Contribution 1: A mathematically grounded TDA pipeline for wearable fall
detection.}  We construct a feature map $\phi_\lambda\colon\R^{n\times 3}\to\R^{24}$
that transforms a raw three-axis accelerometer window into a fixed 24-dimensional
topological descriptor via delay embedding, simplicial filtration, persistent homology
computation, and vectorization.  The map is parameterised by a hyperparameter tuple
$\lambda=(w_\mathrm{sec},m,\tau,r,c,d_\mathrm{met},\varepsilon,q)$ whose components
control the geometric construction.  The Cohen--Steiner stability theorem provides a
formal noise-robustness certificate for the resulting features.

\textbf{Contribution 2: A classifier-agnostic Bayesian hyperparameter optimisation
framework.}  A single Optuna/TPE study~\parencite{optuna} with a classifier-averaged
$F_2$ objective identifies $\lambda^*_K$ for each dataset~$K$ in 200 trials.  The
averaged objective ensures that the recovered representation supports both logistic
regression and support vector machines without being over-fitted to either.  We show
that this optimisation identifies dataset-specific configurations (H2) and that the
variation in $F_2$ across trials---up to 0.18 units on MobiFall---far exceeds the
variation attributable to classifier choice (H3).

\textbf{Contribution 3: A three-protocol subject-aware evaluation framework.}  Three
evaluation protocols are reported jointly for each dataset and classifier: the Naive
(random split, upper bound), the General (leave-one-subject-out, primary), and the
Personal (per-subject split, idealised personalisation reference).  The LOSO protocol
with $R=15$ independent repetitions provides repetition-level confidence intervals via
a $t$-statistic, and the joint reporting of all three protocols makes the
train--test leakage implicit in naive evaluation visible quantitatively.

\textbf{Contribution 4: A statistical characterisation of threshold heterogeneity
as a distributional phenomenon.}  The subject-level sweep gain
$\Delta F_2(s) = F_2(\tau^*(s), s) - F_2(\hat\tau_0(s), s)$ is analysed as a
distribution rather than a scalar average.  We show that the distribution is
right-skewed and dataset-dependent: MobiFall exhibits a mean gain of ${\approx}0.04$--$0.05$
with a coefficient of variation exceeding $50\%$, while SisFall exhibits a mean gain
near zero for most subjects.  This distributional analysis confirms Hypothesis~4
and replaces the anecdotal treatment of personalisation common in prior work.

\textbf{Contribution 5: A deployment pathway to consumer devices.}  The pipeline
is CPU-only, requires no GPU, and executes persistent homology at the window sizes
considered in tens of milliseconds on a desktop-class CPU core---a profile compatible
with ARM Cortex-A smartphone processors and modern smartwatch SoCs.  The threshold
sweep personalisation runs in under one second on device.  We describe a prospective
on-device calibration procedure that requires only a short wearing period and no
server-side computation, and identify one-class anomaly detection from ADL data as
the natural extension that eliminates the cold-start problem entirely.

% ─────────────────────────────────────────────────────────────
\section{Thesis Organisation}
\label{subsec:organisation}

The remainder of the thesis is organised as follows.

\textbf{Chapter~2 (Related Work)} surveys the three families of prior wearable fall
detection methods---threshold-based, statistical-feature, and deep-learning---and
positions the TDA approach with respect to each.  It also reviews prior applications
of persistent homology to time-series classification.

\textbf{Chapter~3 (Methodology)} provides the mathematical foundations.  It
introduces the signal-magnitude-vector transformation and delay embedding, develops
the theory of simplicial complexes, persistent homology, and the Cohen--Steiner
stability theorem, describes the feature vectorization and the Bayesian optimisation
procedure, and defines the three evaluation protocols and the statistical inference
framework used throughout the thesis.

\textbf{Chapter~4 (Experiments)} describes the three benchmark datasets, the
experimental setup, and the numerical results: dataset-specific optimal configurations,
full performance tables for all classifier-protocol combinations, and the threshold
sweep analysis.

\textbf{Chapter~5 (Results and Discussion)} interprets the experimental outcomes
against the four research hypotheses, presents the per-subject threshold gain
distributions, synthesises findings across datasets, discusses the deployment
pathway, and identifies limitations and directions for future work.

\textbf{Chapter~6 (Conclusion)} summarises the contributions at three levels:
mathematical, statistical, and engineering.

% ─────────────────────────────────────────────────────────────
\section*{Notation}

Throughout the thesis, $\R$ denotes the real line, $\R^d$ the $d$-dimensional
Euclidean space, and $\mathbf{1}\{\cdot\}$ the indicator function.  Scalars are
denoted by lowercase letters, vectors by bold lowercase, and sets by calligraphic
uppercase.  The symbol $\|\cdot\|_2$ denotes the Euclidean norm and
$\|\cdot\|_\infty$ the supremum norm.  All probabilistic statements are with
respect to the randomness of the class-balancing seed unless stated otherwise.

